\section{Discussion and Conceptual Positioning}
  \label{sec:discussion}

  The analysis presented in this work proposes a structural reinterpretation of Bell
  inequality violations that avoids introducing nonlocal dynamics, hidden variables,
  or modifications of quantum mechanics.
  It is therefore important to clarify how this approach relates to existing
  interpretations and foundational proposals addressing quantum nonlocality.

  Unlike local hidden-variable models, the present framework does not assume the
  existence of a set of accessible variables that fully determine observable outcomes~\cite{tHooft2016CAI}.
  No completion of the quantum state is postulated, and no additional degrees of
  freedom are introduced.
  The failure of Bell inequalities arises not from incomplete knowledge of underlying
  variables, but from the non-injective nature of the mapping between underlying
  configurations and observable descriptions.

  The present approach is also distinct from superdeterministic models~\cite{Hall2010LocalDeterministic}.
  While superdeterminism denies statistical independence between measurement settings
  and underlying variables, no such assumption is required here.
  Measurement choices remain free and uncorrelated with the underlying configuration space.
  Bell violations arise solely from the structural properties of effective descriptions,
  not from conspiratorial correlations between settings and states.

  Retrocausal interpretations similarly invoke influences propagating backward in time
  to account for nonlocal correlations.
  By contrast, the mechanism discussed here is entirely atemporal at the descriptive
  level.
  No causal influence, forward or backward, is required between distant measurement
  events.
  Correlations emerge from the global consistency constraints imposed by non-factorizable
  effective descriptions, rather than from dynamical signaling.

  The present framework is compatible with standard accounts of quantum contextuality,
  but shifts the emphasis from contextual measurement dependence to descriptive
  non-injectivity.
  Rather than attributing Bell violations to context-dependent properties of observables,
  we identify a more general structural origin: the impossibility of assigning
  independent local outcomes within a single effective description when the underlying
  mapping is many-to-one.

  Finally, we stress that the present work does not advocate a specific ontological
  interpretation of quantum mechanics.
  It neither endorses nor rejects realism, operationalism, or informational approaches.
  Our aim is more modest: to isolate a minimal and sufficient structural condition under
  which Bell-type inequalities fail.
  This condition is independent of any particular ontological commitment and applies
  to any framework in which observable descriptions arise through non-injective mappings.
  Non-injective projection, understood as a finite-capacity effect, provides such a
  condition, clarifying how quantum correlations can arise without abandoning locality
  at the underlying level.

  In physical realizations, this non-injectivity can be understood as a consequence
  of finite descriptive capacity.
  In particular, Planck's constant $h$ sets a fundamental bound on the
  distinguishability of underlying configurations, thereby controlling the effective
  size of the equivalence classes $[\omega]$ that define observable outcomes.

  In this sense, Bell inequality violations may be viewed not as evidence for exotic
  dynamics, but as indicators of intrinsic limitations of finite-capacity effective
  descriptions~\cite{Spekkens2007Epistemic}.
  They reveal the boundaries of applicability of classical probabilistic reasoning,
  rather than the presence of nonlocal causal mechanisms.
